\section{Вступ}
\label{sec:introduction}

Кутові та ламані межі поділу різнорідних матеріалів виникають у багатьох
клейових і багатоматеріальних з'єднаннях. Зокрема, у подвійних косих
з'єднаннях сама геометрія створює внутрішню сингулярну область, що істотно
впливає на розвиток пошкодження~\cite{GacoinEtAl2009}. В ідеалізованій
пружній постановці кутова точка біматеріальної межі може породжувати
степеневе сингулярне поле, показник та кутова структура якого залежать від
геометрії і пружної невідповідності матеріалів. Класичні умови виникнення
та характеристики таких сингулярностей було встановлено для з'єднаних
ізотропних клинів~\cite{BogyWang1971,DempseySinclair1981}; подальший
аналіз показав, що зміна кута може істотно змінювати не лише показник
сингулярності, а й інтенсивність локального поля
\cite{MohammedLiechti2001}. Цей напрям залишається активним: сучасні
роботи розвивають стійкі числові процедури визначення характеристик
сингулярного поля у тривимірних з'єднаннях
\cite{MiyazakiFukudaNoda2025} та маломасштабні співвідношення, що
пов'язують поле неушкодженої кутової точки з параметрами короткої
тріщини, яка з неї виходить~\cite{LyuZhangMai2026}. Отже, для задач
руйнування принциповим є не лише саме існування кутової сингулярності, а
механізм, за яким локальне пружне поле переходить у зародження та розвиток
тріщини.

Для початково неушкоджених біматеріальних кутів такий перехід
досліджувався як за допомогою локальних параметрів сингулярного поля, так
і в моделях зі скінченною зоною руйнування. Mohammed і Liechti поєднали
експериментальні вимірювання з когезійним моделюванням зародження тріщини
для різних кутів біматеріального з'єднання~\cite{MohammedLiechti2000}.
Labossiere і Dunn поширили підхід, заснований на інтенсивностях
сингулярних полів, на тривимірні біматеріальні кути
\cite{LabossiereDunn2001}. Akisanya і Meng показали, що використання
критичного параметра пружного кутового поля як характеристики руйнування
потребує достатнього розділення масштабів між нелінійною зоною та
областю, в якій домінує асимптотичне сингулярне поле
\cite{AkisanyaMeng2003}. У сучасних когезійних формулюваннях роль
характерного масштабу зони руйнування сформульована ще виразніше:
можливість наближено трактувати кут як тріщиноподібний концентратор
залежить від відношення когезійної довжини до геометричного масштабу
задачі~\cite{GormanThouless2022,ThoulessGoutianos2023}. Ці результати
безпосередньо обґрунтовують маломасштабний характер зони
передруйнування: її малий розмір відносно зовнішнього масштабу є не лише
технічною асимптотичною передумовою, а й частиною фізичної області
застосовності моделі.

Коли міжфазна тріщина вже існує, виникає інша задача --- конкуренція між
подальшим поширенням уздовж межі та виходом тріщини в один із матеріалів.
Класичні енергетичні дослідження сформулювали умови початкового
відхилення тріщини від межі~\cite{HeHutchinson1989}, а для анізотропних
біматеріалів було отримано співвідношення для коефіцієнтів інтенсивності
та швидкостей вивільнення енергії при такому відхиленні
\cite{Wang1994}. Когезійні моделі уточнюють цю картину: вибір шляху може
залежати не лише від тріщиностійкості, а й від когезійної міцності та
характерної довжини зони руйнування~\cite{ParmigianiThouless2006}.
Аналітичні моделі зон передруйнування безпосередньо описують малу зону,
що виходить із вершини плоскої міжфазної тріщини в один із матеріалів, і
дають її довжину, напрям, розкриття та енергетичні характеристики
\cite{KaminskyEtAl2023}. Споріднена модель для квазікрихкого
з'єднувального матеріалу описує маломасштабну процес-зону при подальшому
міжфазному розвитку тріщини~\cite{KaminskyDudykChornoivan2025}, а
урахування несингулярного $T$-напруження показує, що в задачах
відгалуження вищі члени зовнішнього поля можуть давати помітні кількісні
поправки~\cite{KaminskyDudykChornoivan2026}. Ці результати формують методичну основу для подальшого розвитку
процес-зонного підходу. Специфіка цієї роботи визначається вихідною
конфігурацією контактуючих міжфазних зсувних тріщин у кутовій точці
ламаної межі та механічним переходом до відгалуження.

До розглядуваної тут геометрії найближчими є дослідження тріщин і зон
локального руйнування безпосередньо біля кутової точки ламаної межі.
Початкову стадію відхилення міжфазної тріщини в кутовій точці було
проаналізовано Dudyk і Dikhtyarenko~\cite{DudykDikhtyarenko2012}, а
маломасштабну зону руйнування, що ініціюється безпосередньо в кутовій
точці V-подібної межі, --- Kaminsky, Kipnis і Polishchuk
\cite{KaminskyKipnisPolishchuk2012}. Для тріщини, що виходить із
кутової точки ламаної межі, Kaminsky, Dudyk і Reshitnyk побудували
аналітичну модель когезійної зони передруйнування з використанням
перетворення Мелліна та методу Вінера--Гопфа
\cite{KaminskyDudykReshitnyk2020}. Для міжфазної тріщини з вершиною в
кутовій точці полігональної межі параметри малої зони руйнування було
визначено також за контактної взаємодії берегів
\cite{KaminskyDudykFenkiv2023}. Окрема близька лінія досліджень охоплює
пластичну зону біля контактуючої міжфазної тріщини
\cite{KaminskyDudyk2024PlasticContact}, квазікрихку процес-зону біля
міжфазної тріщини в кутовій точці ламаної межі з енергетичним вибором її
напряму~\cite{BogdanovaEtAl2024} та ініціювання тріщини біля ламаної
межі за наявності пластичного з'єднувального шару
\cite{KaminskyDudykPolishchuk2024}. Отже, процес-зони біля тріщин у
кутових точках і сам вихід тріщини з межі поділу вже мають розвинену
теоретичну основу. Відмінність теперішньої задачі полягає в іншому
вихідному стані та в механічній передісторії відгалуження.

Безпосередньою вихідною точкою цієї роботи є отримане Nazarenko і Kipnis
сингулярне поле біля кутової точки ламаної межі поділу за наявності
симетричної міжфазної тріщини зсуву з контактуючими без тертя берегами
\cite{NazarenkoKipnis2022Stress}. У цій роботі дві симетричні гілки
розглядаються як одна V-подібна міжфазна тріщина зсуву: кожна її гілка
має довжину $l$, а сумарна довжина гілок дорівнює $2l$. Для тієї самої
конфігурації в іншій роботі Nazarenko і Kipnis було визначено коефіцієнт
інтенсивності напружень за модою~II біля зовнішніх вершин тріщини,
сформульовано умову її граничної рівноваги та показано нестійкість цієї
рівноваги щодо подальшого міжфазного поширення
\cite{NazarenkoKipnis2022Equilibrium}. Цей механізм локалізований біля
зовнішніх вершин і не є вихідною стадією розглядуваної тут локальної
задачі. V-подібна міжфазна тріщина вважається заданою стаціонарною
конфігурацією, а рівень зовнішнього навантаження --- недостатнім для
втрати нею локальної стійкості.

Предметом цієї роботи є локальний механізм ініціювання відгалуження
V-подібної міжфазної тріщини зсуву в її кутовій точці, пов'язаний з
утворенням малої розтягувальної зони передруйнування в одному з
матеріалів. За маломасштабного наближення довжину цієї зони в матеріалі
$i$ ($i=1,2$) позначаємо через $d_i$ і вважаємо значно меншою за довжину
гілки тріщини $l$. Поле V-подібної тріщини без зони використовується як
зовнішнє для локальної задачі біля кутової точки, тоді як зворотний вплив
малої зони на зовнішні вершини тріщини у провідному наближенні не
враховується. Необхідно встановити, у якому з матеріалів така зона є
фізично допустимою, та визначити її рівноважну довжину $d_i$, розкриття
$\delta_i$, роботу сил опору відриву $W_i$, енергетичний параметр
$J_i=dW_i/dd_i$ і зміну локальної кутової сингулярності.

Метою роботи є побудова аналітичної моделі квазістатичних станів зони
передруйнування, яка дає змогу описати ініціювання відгалуження
V-подібної міжфазної тріщини зсуву в її кутовій точці. В умовах плоскої
деформації та симетричного навантаження зона моделюється лінією розриву
нормального переміщення, орієнтованою вздовж бісектриси кутового сектора
відповідного матеріалу, на якій діє стале нормальне напруження опору
відриву. За допомогою перетворення Мелліна, зрощування асимптотик і
методу Вінера--Гопфа визначаються умови фізичної допустимості,
рівноважна довжина та розкриття зони, робота сил опору відриву,
пов'язаний із зоною енергетичний параметр і показник залишкової кутової
сингулярності. Розкриття та енергетичний параметр можуть бути використані
відповідно в деформаційному й енергетичному критеріях ініціювання
відгалуженої тріщини після незалежного задання критичної характеристики
матеріалу. Подальше поширення та траєкторія скінченного відгалуження до
цієї локальної постановки не входять.
